% 2015 全国Ⅱ文科 18(1) 解答图：B 地区频率分布直方图。
% 对照源卷第 968 页空白网格；横轴保留 40 左端界及 50--100 刻度，纵轴至 0.040。
\begin{tikzpicture}[font=\small, x=0.105cm, y=125cm]
  % 原卷网格（仅绘图区域，不复录题干）。
  \foreach \x in {40,50,...,100}
    \draw[gray!65, line width=0.25pt] (\x,0) -- (\x,0.04);
  \foreach \y in {0.005,0.010,...,0.040}
    \draw[gray!65, line width=0.25pt] (40,\y) -- (100,\y);

  % 五组柱高为频率/组距：0.005、0.020、0.035、0.025、0.015。
  \foreach \left/\right/\height in {
      50/60/0.005,
      60/70/0.020,
      70/80/0.035,
      80/90/0.025,
      90/100/0.015}
    \filldraw[fill=gray!18, draw=black, line width=0.55pt]
      (\left,0) rectangle (\right,\height);

  % 坐标轴与刻度。
  \draw[-stealth, line width=0.65pt] (40,0) -- (103,0);
  \draw[-stealth, line width=0.65pt] (40,0) -- (40,0.0435);
  \foreach \x in {50,60,...,100} {
    \draw (\x,0) -- (\x,-0.00065);
    \node[below, inner sep=2pt] at (\x,-0.00065) {\footnotesize \x};
  }
  \foreach \y/\label in {
      0.005/0.005,
      0.010/0.010,
      0.015/0.015,
      0.020/0.020,
      0.025/0.025,
      0.030/0.030,
      0.035/0.035,
      0.040/0.040} {
    \draw (40,\y) -- (39.2,\y);
    \node[left, inner sep=2pt] at (39.2,\y) {\footnotesize \label};
  }
  \node[below left, inner sep=1.5pt] at (40,0) {$O$};
  \node[right, inner sep=3pt] at (103,0) {满意度评分};
  \node[above, inner sep=2pt] at (40,0.0435) {频率/组距};
\end{tikzpicture}
